\subsection{Gesamtarchitektur}
\label{subsec:gesamtarchitektur}

Die Baseline besteht aus einem Master-Repository mit gemeinsamem Kern und
deklarativen Varianten. Profile wählen wiederverwendbare Features; optionale
Capabilities erweitern nur kompatible Profile. Diese beherrschte Variabilität
vermeidet getrennte Template-Kopien
\parencite{clements2001software-product-lines,kleiveist2026profiles}.

\begin{figure}[H]
  \centering
  \resizebox{0.96\textwidth}{!}{%
  \begin{tikzpicture}[node distance=7mm and 10mm]
    \node[casePrimary,minimum width=48mm] (master)
      {Master-Repository\\vollständige Baseline};
    \node[caseBox,minimum width=40mm,below left=of master] (catalog)
      {\textbf{profiles/}\\Features und Presets};
    \node[caseBox,minimum width=40mm,below right=of master] (selection)
      {Profilwahl\\optionale Capabilities};
    \node[casePrimary,minimum width=48mm,below=11mm of master] (generator)
      {Generator\\Auswahl und Abhängigkeiten auflösen};
    \node[caseBox,minimum width=48mm,below=of generator] (project)
      {abgeleitetes Projekt\\\texttt{project-profile.toml}};

    \node[caseBox,minimum width=30mm,below left=10mm and 52mm of project] (frontend)
      {Frontend};
    \node[caseOptional,minimum width=30mm,right=of frontend] (backend)
      {Backend};
    \node[caseOptional,minimum width=30mm,right=of backend] (tauri)
      {Tauri};
    \node[caseOptional,minimum width=30mm,right=of tauri] (database)
      {Persistenz};

    \draw[caseFlow] (master) -- (catalog);
    \draw[caseFlow] (master) -- (selection);
    \draw[caseFlow] (catalog) -- (generator);
    \draw[caseFlow] (selection) -- (generator);
    \draw[caseFlow] (generator) -- (project);
    \draw[caseFlow] (project) -- (frontend);
    \draw[caseSecondaryFlow] (project) -- (backend);
    \draw[caseSecondaryFlow] (project) -- (tauri);
    \draw[caseSecondaryFlow] (project) -- (database);

    \begin{scope}[on background layer]
      \node[caseGroup,fit=(frontend)(database),
        label={[font=\sffamily\scriptsize,text=black!60]below:nur aufgelöste Features werden übernommen}] {};
    \end{scope}
  \end{tikzpicture}}
  \caption{Gesamtarchitektur des Technologie-Templates}
  \label{fig:system-architecture}
  \smallskip
  {\footnotesize Quelle: Eigene Darstellung auf Grundlage von
  \textcite{kleiveist2026architecture} und \textcite{kleiveist2026profiles}.\par}
\end{figure}


Abbildung~\ref{fig:system-architecture} trennt die Baseline vom abgeleiteten
Projekt. Der Generator löst Profil und Capabilities auf, übernimmt die
zugehörigen Pfade und dokumentiert das Ergebnis in
\texttt{project-profile.toml}. Frontend, Backend, Tauri und die serverseitige
SQL-Persistenz bleiben eigenständige, teils optionale Bausteine
\parencite{kleiveist2026architecture,kleiveist2026persistence}.

\begin{figure}[H]
  \centering
  \resizebox{0.97\textwidth}{!}{%
  \begin{tikzpicture}[node distance=5.5mm and 11mm]
    \node[casePrimary,minimum width=48mm] (repository) {Master-Repository};
    \node[font=\sffamily\bfseries\scriptsize,text=caseBlue,above=5mm of repository]
      {Verantwortungsbereiche im Repository};

    \node[casePrimary,minimum width=46mm,below left=10mm and 53mm of repository] (runtime)
      {Laufzeitquellen};
    \node[casePrimary,minimum width=46mm,below=10mm of repository] (contracts)
      {Auswahl und Verträge};
    \node[casePrimary,minimum width=46mm,below right=10mm and 53mm of repository] (lifecycle)
      {Lebenszyklus};

    \node[caseBox,minimum width=46mm,below=of runtime] (frontend)
      {\textbf{frontend/}\\UI, Build, HTTP-Client};
    \node[caseBox,minimum width=46mm,below=of frontend] (backend)
      {\textbf{backend/}\\API, Konfiguration, optionale DB};
    \node[caseBox,minimum width=46mm,below=of backend] (tauri)
      {\textbf{src-tauri/}\\native Desktop-Grenze};

    \node[caseBox,minimum width=46mm,below=of contracts] (profiles)
      {\textbf{profiles/}\\Features und Presets};
    \node[caseBox,minimum width=46mm,below=of profiles] (config)
      {\textbf{config/}\\Variablenvertrag};
    \node[caseBox,minimum width=46mm,below=of config] (shared)
      {\textbf{shared/}\\Schema, Beispiele, Assets};

    \node[caseBox,minimum width=46mm,below=of lifecycle] (tools)
      {\textbf{tools/}\\Generierung und Steuerung};
    \node[caseBox,minimum width=46mm,below=of tools] (deployment)
      {\textbf{deployment/}\\providerneutrale Artefakte};
    \node[caseBox,minimum width=46mm,below=of deployment] (docs)
      {\textbf{docs/}\\Architektur und Betrieb};

    \draw[caseFlow] (repository) -- (runtime);
    \draw[caseFlow] (repository) -- (contracts);
    \draw[caseFlow] (repository) -- (lifecycle);

    \begin{scope}[on background layer]
      \node[caseGroup,fit=(runtime)(frontend)(backend)(tauri)] {};
      \node[caseGroup,fit=(contracts)(profiles)(config)(shared)] {};
      \node[caseGroup,fit=(lifecycle)(tools)(deployment)(docs)] {};
    \end{scope}
  \end{tikzpicture}}
  \caption{Komponenten und Verantwortungsgrenzen im Master Repository}
  \label{fig:component-architecture}
  \smallskip
  {\footnotesize Quelle: Eigene Darstellung auf Grundlage von
  \textcite{kleiveist2026architecture}.\par}
\end{figure}


Die Komponententrennung in Abbildung~\ref{fig:component-architecture} hält
Client, Server, native Integration und gemeinsame Artefakte auseinander.
Profile und Konfiguration beschreiben Struktur und Laufzeitwerte; Tooling und
Deployment steuern die Komponenten außerhalb der Produktlaufzeit. Seit
\texttt{461fc751} gehört die Quality-Governance zu diesem Entwicklungs- und
CI-Bereich. Sie untersucht Laufzeitquellen, wird aber weder ausgeliefert noch
Teil des HTTP-, Tauri- oder Persistenzdatenflusses
\parencite{kleiveist2026governance-snapshot}.

Lokale Entwickler, Coding Agents und GitHub Actions rufen
\path{tools/control.py} auf. Der Dispatcher delegiert an die getrennte
Quality-Orchestrierung; diese verbindet zentrale Policy, Repository-Scan,
sprachspezifische Werkzeuge, Backend- und Frontend-Architekturprüfungen,
Exception-Behandlung sowie Reporting. Abbildung~\ref{fig:verification-model}
stellt diesen Kontrollfluss in Abschnitt~\ref{subsec:teststrategie} dar. Die
Architekturentscheidung ergänzt die vorhandene Laufzeitarchitektur um eine
Prüfschicht, nicht um eine weitere Produktkomponente.

Der Governance-Commit korrigierte zugleich die zuvor direkte Kopplung des
Health-Routers an SQLAlchemy und \path{app.db}: Der Router delegiert nun an
einen Service; konkrete Infrastruktur wird in \path{backend/app/main.py}
verdrahtet. Dies bildet einen beobachtbaren Vorher-Nachher-Fall der später
automatisierten Backend-Grenzen \parencite{kleiveist2026governance-commit}.

\begin{figure}[H]
  \centering
  \resizebox{0.88\textwidth}{!}{%
  \begin{tikzpicture}[node distance=8mm and 34mm]
    \node[casePrimary,minimum width=34mm] (browser) {Browser};
    \node[casePrimary,minimum width=34mm,right=of browser] (desktop) {Tauri-Webview};
    \node[caseBox,minimum width=70mm,below=17mm of $(browser)!0.5!(desktop)$] (frontend)
      {ein gemeinsamer Vite-/TypeScript-Build};
    \node[caseOptional,minimum width=48mm,below=of frontend] (api)
      {FastAPI-Service\\nur backendfähige Profile};
    \node[caseOptional,minimum width=48mm,below=of api] (postgres)
      {SQLAlchemy + Psycopg\\PostgreSQL optional};

    \draw[caseFlow] (browser) -- node[pos=0.42,left,font=\sffamily\tiny] {lädt} (frontend);
    \draw[caseFlow] (desktop) -- node[pos=0.42,right,font=\sffamily\tiny] {hostet} (frontend);
    \draw[caseSecondaryFlow] (frontend) -- node[right,font=\sffamily\tiny] {HTTP /api} (api);
    \draw[caseSecondaryFlow] (api) -- node[right,font=\sffamily\tiny] {serverseitig} (postgres);

    \begin{scope}[on background layer]
      \node[caseGroup,fit=(browser)(desktop)(postgres),
        label={[font=\sffamily\scriptsize,text=black!60]below:Komponentenbestand abhängig vom aufgelösten Profil}] {};
    \end{scope}
  \end{tikzpicture}}
  \caption{Profilabhängige Laufzeitarchitektur}
  \label{fig:runtime-architecture}
  \smallskip
  {\footnotesize Quelle: Eigene Darstellung auf Grundlage von
  \textcite{kleiveist2026architecture} und \textcite{kleiveist2026profiles}.\par}
\end{figure}


Browser oder Tauri-Webview laden denselben Vite-Build
(Abbildung~\ref{fig:runtime-architecture}). Backendfähige Profile nutzen
FastAPI über HTTP; allein das Backend greift auf die dargestellte optionale
PostgreSQL-Persistenz zu. Lokale Produktspeicher sind davon getrennt und nicht
implementiert \parencite{kleiveist2026persistence}.
\beginnerinput{workspace/statement/300_architektur/310}
